% cSpell:words DTFT Vieta
\section{Homework 12}

\subsection{序列的傅里叶变换}

(1) \begin{equation*}
    DTFT[y[n]]=\left.X(z)\right|_{z=e^{j\omega}}=\frac{1}{1-\frac{1}{2}e^{-j\omega}}
\end{equation*}

(2) \begin{equation*}
    DTFT[y[n]]=\left.X(z)\right|_{z=e^{j\omega}}=\frac{e^{j\omega}(e^{j\omega}+\frac{1}{2})}{e^{2j\omega}-e^{j\omega}+1}
\end{equation*}

(3) \begin{equation*}
    DTFT[y[n]]=\left.X(z)\right|_{z=e^{j\omega}}=1+e^{-j\omega}+e^{-2j\omega}+e^{-3j\omega}+e^{-4j\omega}
\end{equation*}

\subsection{系统框图与系统函数}

\subsubsection{连续时间系统}

\task{第一小题(必做题)}

(1) 由系统框图可列出: $$\frac{1}{s+2}\left(\frac{1}{s+1}\left(X(s)-2Y(s)\right)-2Y(s)\right)=Y(S),$$
解得$H(s)=\dfrac{1}{s^2+5s+6}$

(2) $H(s)=\dfrac{1}{s^2+5s+6}=\dfrac{1}{s+2}-\dfrac{1}{s+3}$, 因此$h(t)=\mathscr{L}^{-1}[H(s)]=[e^{-2t}+-e^{-3t}]u(t)$.

(3) 由于$h(t)$积分绝对收敛, 该系统是稳定系统.

\subsubsection{离散时间系统}

\task{第一小题(必做题)}

(1) 设中间节点的信号为$A(z)$, 则由系统框图可列方程\begin{equation*}
    \left\{
        \begin{array}{ll}
            X(z)-4z^{-1}A(z)=A(z) \\
            2A(z)+z^{-1}A(z)=Y(z)
        \end{array}
    \right.
\end{equation*}
可解出$H(z)=\frac{2+z^{-1}}{1+4z^{-1}}$.

(2) $X(z)=\dfrac{1}{1+z^{-1}}+\dfrac{1}{1-\frac{1}{2}z^{-1}}$, 因此$Y(z)=X(z)H(z)=\dfrac{2+z^{-1}}{(1+z^{-1})(1-\frac{1}{2}z^{-1})(1+4z^{-1})}\\=-\dfrac{2}{9}\cdot\dfrac{1}{1+z^{-1}}+\dfrac{4}{27}\cdot\dfrac{1}{1-\frac{1}{2}z^{-1}}+\dfrac{56}{27}\cdot\dfrac{1}{1+4z^{-1}}$, 因此$y_{zs}[n]=-\dfrac{2}{9}\cdot(-1)^n+\dfrac{4}{27}\cdot\left(\dfrac{1}{2}\right)^n+\dfrac{56}{27}\cdot(-4)^n$.

(3) 因为$H(z)=\dfrac{1}{4}+\dfrac{7}{4}\cdot\dfrac{1}{1+4z^{-1}}$, 有$h[n]=\dfrac{1}{4}\delta[n]+\dfrac{7}{4}(-4)^nu[n]$, 只在$\ge0$处取值且级数和绝对收敛, 因此该系统是因果稳定系统.

\task{第二小题(必做题)}

(1) 设$H(z)=\dfrac{K}{(1-\frac{1}{3}z^{-1})(1-\frac{2}{5}z^{-1})}=K\left(-5\cdot\dfrac{1}{1-\frac{1}{3}z^{-1}}+6\cdot\dfrac{1}{1-\frac{2}{5}z^{-1}}\right)$, 因此$h[n]=K\left(-5\cdot\left(\dfrac{1}{3}\right)^n+6\cdot\left(\dfrac{2}{5}\right)^n\right)$. 代入初值$h[0]=3$, 得到$K=3$.

(2) $h[n]=3\left(-5\cdot\left(\dfrac{1}{3}\right)^n+6\cdot\left(\dfrac{2}{5}\right)^n\right)$.

\subsection{利用系统函数分析输入输出}

(1) $X(s)=\dfrac{1}{s+1}$, $Y(s)=\dfrac{1}{2}\dfrac{1}{s+1}-2\cdot\dfrac{1}{s+2}+\dfrac{1}{s+3}$, 因此$H(s)=Y(s)/X(s)=-\dfrac{1}{2}-\dfrac{2(s+1)}{s+2}+\dfrac{s+1}{s+3}=-\dfrac{1}{2}+\dfrac{2}{s+2}-\dfrac{2}{s+3}$, 因此$h(t)=-\dfrac{1}{2}\delta(t)+2e^{-2t}-2e^{-3t}$.

(2) $H(s)=1+\frac{1}{s+3}$, $Y(s)=\dfrac{2}{s+2}-\dfrac{1}{s+3}$, 因此$X(s)=Y(s)/H(s)=\dfrac{2(s+3)}{s+2}-1$, $x(t)=\delta(t)+2e^{-2}$.

(3) 对此差分方程两边做Z变换: $Y(z)-\dfrac{1}{2}z^{-1}Y(z)=X(z)$, 代入$Y_{zs}(z)=\dfrac{1}{1-\frac{1}{2}z^{-1}}-\dfrac{1}{1-\frac{1}{3}z^{-1}}$, 得到$X(z)=-\dfrac{1}{2}+\dfrac{1}{2}\cdot\dfrac{1}{1-\frac{1}{3}z^{-1}}$, 因此$x[n]=-\dfrac{1}{2}\delta[n]+\dfrac{1}{2}\dfrac{1}{3^n}$.

\subsection{系统的稳定性和因果性}

\task{第一小题(必做题)}

(1) $\dfrac{1}{s^2+5s-6}(X(s)-kY(s))=Y(s)\Rightarrow H(s)=\dfrac{1}{s^2+5s-6+k}$.

(2) 要求所有的极点都在左半侧, 因此由Vieta定理, 有$-6+k>0$, 得到$k>6$.

(3) $s^2+5s-6+k=(s+\frac{5}{2})^2-\frac{49}{4}+k$, 因此$k=\frac{49}{4}$, 此时$H(s)=\dfrac{1}{(s+\frac{5}{2})^2}$, $h(t)=te^{-\frac{5}{2}t}$.

\task{第三小题(必做题)}

(1) $Y(z)-\dfrac{9}{10}z^{-1}Y(z)=X(z)-\dfrac{1}{2}z^{-1}X(z)$, 因此$H(s)=\dfrac{1-\frac{1}{2}z^{-1}}{1-\frac{9}{10}z^{-1}}$, 有极点$z=\dfrac{10}{9}>1$, 因此不为稳定系统

(2) $Y(z)+2z^{-1}Y(z)+z^{-2}Y(z)=\dfrac{1}{2}X(z)-z^{-2}X(z)$, 有$H(z)=\dfrac{\frac{1}{2}-z^{-2}}{1+2z^{-1}+z^{-2}}$, 没有极点, 是稳定系统.

\task{第四小题(必做题)}

对于Z变换表达式, 因果代表收敛域在无穷远处有定义. 对于(2), (3), 分子的最高指数大于分母的最高指数, 因此在无穷远处没有意义, 不为因果信号, 对于(1), 无穷远处的极限为$1$, 是因果系统.

\subsection{电路系统分析}

设$R_1$, $R_2$中间节点的电位为$v_4(t)$, 则节点电压法列出方程:\begin{equation*}
    \left\{
        \begin{array}{ll}
            \frac{v_1-v_4}{R_1}+C_2\frac{\mathrm d(v_3-v_4)}{\mathrm dt}-C_1\frac{\mathrm dv_2}{\mathrm dt}=0 \\
            v_2-R_2C_1\frac{\mathrm dv_2}{\mathrm dt}=v_4 \\
            v_3=kv_2
        \end{array}
    \right.
\end{equation*}

解出$\dfrac{1}{k^2}\dfrac{\mathrm d^2v_3}{\mathrm dt^2}+(1-\dfrac{1}{k})\dfrac{\mathrm dv_3}{\mathrm dt}-\dfrac{1}{k}v_3=-v_1$, 因此系统函数$$V_3(s)/V_1(s)=-\dfrac{1}{\frac{1}{k^2}s^2+(1-\frac{1}{k})s-\frac{1}{k}},$$

令$k=2$, 得到$H(s)=-\dfrac{4}{s^2+2s-2}=-\dfrac{2}{\sqrt{3}}(\dfrac{1}{s+1-\sqrt{3}}-\dfrac{1}{s+1+\sqrt{3}})$, $h(t)=\\-\dfrac{2}{\sqrt{3}}(e^{(-1+\sqrt{3})t}+e^{(-1-\sqrt{3})t})u(t)$, 积分得到单位阶跃响应$y(t)=\dfrac{2}{3-\sqrt{3}}e^{(-1+\sqrt{3})t}+\\\dfrac{2}{3+\sqrt{3}}e^{(-1-\sqrt{3})t}$ 

